\subsection{Taylor展開と数値微分}

\subsubsection{Taylor展開}

$f$ が $x=a$ の近傍で十分な回数だけ微分可能であるとする。
\[
 f(x)=C_0+C_1(x-a)+C_2(x-a)^2+\cdots
\]
とおき、両辺を繰り返し微分して $x=a$ を代入すると
$C_n=f^{(n)}(a)/n!$ を得る。したがって
\begin{equation}
 \boxed{f(x)=\sum_{n=0}^{\infty}
 \frac{f^{(n)}(a)}{n!}(x-a)^n}
\end{equation}
を $a$ のまわりの\textbf{Taylor展開}という。有限次数で打ち切るときは
\[
 f(x)=\sum_{k=0}^{n}\frac{f^{(k)}(a)}{k!}(x-a)^k+R_n(x)
\]
と書き、$R_n$ を剰余項という。Taylorの定理によれば、適切な
$\xi$ が $a$ と $x$ の間に存在して
\[
 R_n(x)=\frac{f^{(n+1)}(\xi)}{(n+1)!}(x-a)^{n+1}
\]
と表せる。

$a=0$ の場合を\textbf{Maclaurin展開}という。代表例は
\begin{align}
 e^x&=1+x+\frac{x^2}{2!}+\frac{x^3}{3!}+\cdots,\\
 \sin x&=x-\frac{x^3}{3!}+\frac{x^5}{5!}-\cdots,\\
 \cos x&=1-\frac{x^2}{2!}+\frac{x^4}{4!}-\cdots,\\
 (1+x)^\alpha
 &=1+\alpha x+\frac{\alpha(\alpha-1)}{2!}x^2+\cdots
 \qquad(|x|<1)
\end{align}
である。たとえば
\[
 \sqrt{1+x}=1+\frac12x-\frac18x^2+O(x^3),
\qquad
 \frac1{1+x}=1-x+x^2+O(x^3).
\]
ここで $O(x^m)$ は、$x\to0$ のとき絶対値が定数倍の $|x|^m$
以下で抑えられる項を表す。

\paragraph{係数計算の例}
$f(x)=\sin x$ では
\[
 f(0)=0,\quad f'(0)=1,\quad f''(0)=0,\quad
 f^{(3)}(0)=-1,\quad f^{(4)}(0)=0
\]
なので、3次までなら
\[
 \sin x=x-\frac{x^3}{3!}+O(x^5).
\]
$f(x)=\sqrt{1+x}$ では
\[
 f(0)=1,\qquad f'(x)=\frac12(1+x)^{-1/2},\qquad
 f''(x)=-\frac14(1+x)^{-3/2}
\]
より
\[
 f'(0)=\frac12,\qquad f''(0)=-\frac14,
\]
したがって
\[
 \sqrt{1+x}
 =1+\frac12x+\frac{f''(0)}{2!}x^2+O(x^3)
 =1+\frac12x-\frac18x^2+O(x^3)
\]
となる。

\subsubsection{数値微分}

関数値が離散的な点でしか得られない場合、導関数を差分で近似する。
Taylor展開
\begin{align*}
 f(x+h)&=f(x)+f'(x)h+\frac12f''(x)h^2+O(h^3),\\
 f(x-h)&=f(x)-f'(x)h+\frac12f''(x)h^2+O(h^3)
\end{align*}
から、前進差分と後退差分
\begin{align}
 f'(x)&=\frac{f(x+h)-f(x)}h+O(h),\\
 f'(x)&=\frac{f(x)-f(x-h)}h+O(h)
\end{align}
を得る。両者を平均すると1次の誤差が打ち消され、
\begin{equation}
 \boxed{f'(x)=\frac{f(x+h)-f(x-h)}{2h}+O(h^2)}
\end{equation}
となる。これを\textbf{中心差分}という。

\begin{center}
\begin{tikzpicture}[>=Stealth,scale=1.1]
 \draw[->] (-.4,0)--(5,0) node[right] {$x$};
 \draw[->] (0,-.2)--(0,3.5) node[above] {$y$};
 \draw[thick,domain=.4:4.5,samples=70,smooth]
   plot(\x,{.18*(\x-1.6)^2+.55});
 \foreach \x/\lab in {1.5/{x-h},2.5/x,3.5/{x+h}}
  {\draw[dashed] (\x,0)--(\x,{.18*(\x-1.6)^2+.55});
   \node[below] at (\x,0) {$\lab$};}
 \draw[blue] (2.5,.45)--(4.5,1.8) node[right] {前進};
 \draw[red] (.5,.45)--(2.5,.72) node[left] {後退};
 \draw[green!60!black,thick] (1.1,.2)--(4,1.42) node[right] {中心};
\end{tikzpicture}

\small 前進差分・後退差分・中心差分による傾きの近似
\end{center}

2階微分は二つのTaylor展開を加えることで
\begin{align*}
 f(x+h)+f(x-h)
 &=2f(x)+f''(x)h^2\\
 &\quad+\frac1{12}f^{(4)}(x)h^4+O(h^6)
\end{align*}
となるため、
\begin{equation}
 \boxed{f''(x)=\frac{f(x+h)-2f(x)+f(x-h)}{h^2}+O(h^2)}
\end{equation}
と近似できる。$h$ を小さくすれば打切り誤差は減少するが、極端に
小さくすると浮動小数点の丸め誤差や桁落ちが増えるため、実計算では
両者の釣合いを考える必要がある。
